% Journal Article
% LaTeX Template
% Version 2.0 (February 7, 2023)
%
% This template originates from:
% https://www.LaTeXTemplates.com
%
% Author:
% Vel (vel@latextemplates.com)
%
% License:
% CC BY-NC-SA 4.0 (https://creativecommons.org/licenses/by-nc-sa/4.0/)
%
% NOTE: The bibliography needs to be compiled using the biber engine.
%
%%%%%%%%%%%%%%%%%%%%%%%%%%%%%%%%%%%%%%%%%

%----------------------------------------------------------------------------------------
%	PACKAGES AND OTHER DOCUMENT CONFIGURATIONS
%----------------------------------------------------------------------------------------



\documentclass{article}

\usepackage{booktabs} % Required for better horizontal rules in tables

\usepackage{array} % Required for manipulating table columns
\usepackage{caption} % Required for customizing captions
\usepackage{enumitem} % Required for list customization
\usepackage{hyperref} % Required for links
\usepackage{graphicx} % Required for including images

\newcommand{\important}[1]{\begin{center}\fbox{
    \vspace{0.1cm}
    \begin{minipage}{0.9\textwidth}
    #1
    \end{minipage}
    \vspace{0.1cm}
    }\end{center}
    }



\title{Journal} % Article title, use manual lines breaks (\\) to beautify the layout


\author{Nathan Berrit}


\begin{document}

\maketitle % Output the title section



\section{Reviews}

\subsection*{Optical absorption and scattering spectroscopies of single nano-objects (2014)}
Questions :


\subsubsection*{I : Introduction}

\subsubsection*{II : Single-particle detection and spectroscopy: absorption and scattering}

à la fin, on dit qu'il y a des fabry perot effect : est-ce que les anneaux qu'on voit sont les anneaux d'un fabry pérot créé par le substrat ?

\begin{itemize}
    \item Read a paper on SNOM techniques, parce que j'ai rien compris.
\end{itemize}


\section{Articles}

\subsection*{Direct Measurement of the Single-Metal-Cluster Optical Absorption (2004)}

Auteurs : Del Fatti, Vallée, A. Arbouet.

\underline{\textbf{Avantages de la SMS}} :
\begin{itemize}
    \item La ref 10 parle de la modulation de frequence, approche analogue de la SMS
    \item Avantage de la SMS : permet de descendre en puissance incidente à des $\mu W$, alors que les techniques en interférence demandaient des mW.
\end{itemize}

\underline{\textbf{Optimisation du signal SMS}} :
\begin{itemize}
    \item Il faut un petit faisceau pour I'' grand.
    \item optimisation de dy : val optimale pour dy = FWHM pour 2f et dy = 0.8 FWHM pour f. A tester numériquement.
    \item dT/T must then be calculated numerically !
\end{itemize}

\underline{\textbf{System studied}} :
\begin{itemize}
    \item Gold nanosphere with diameters from 5 to 0 nm
    \item substrate : glass 150 $\mu$m thick.
\end{itemize}

\underline{\textbf{Experimental setup}} :
\begin{itemize}
    \item 50x50$\mu$m piezo modulated as $\sim$1 kHz.
    \item 2 lasers : ND:YVO4 at 532 nm or HeNe at 633 nm.
    \item objective to focus at diff limit.
\end{itemize}

\underline{\textbf{Results}} : 
\begin{itemize}
    \item From maesurmements they deduced $\sigma_{ext}$ for a lot of nano-objects and compared it to expected theorical value with good agreement.
    \item The dispersion on dT/T was in agreement with the dispersion on size measured with TEM.
    \item The dT/T goes with theory when we add a surface term on permittivity.
\end{itemize}


\subsection*{Electron and Lattice Heating Contributions to the Transient Optical
Response of a Single Plasmonic Nano-Object}

Auteur : Romain Rouxelle.



\subsection*{High-Pressure Effect on the Optical Extinction
of a Single Gold Nanoparticle (2018)}
\end{document}